\section{Related Work}
\label{sec:related}

This section positions GLASS within earlier work. It first reviews graph-based intrusion detection and language models for security data. It then examines how structural and semantic representations are fused, reviews graph--language feedback mechanisms, and finally explains how this study relates to that work.

\subsection{Graph-based network intrusion detection}

Intrusion detection is conventionally divided by detection methodology. The guidance this
report follows names three primary classes, signature-based, anomaly-based and stateful
protocol analysis \citep{scarfone2007idps}, and the first two can both be posed over a flow
graph. A flow
graph makes the edge the natural prediction unit when the object to be classified is a
communication event between two endpoints. Nodes represent communicating entities such as
IP addresses, and edges represent flows or aggregated flows. This formulation preserves the
relational context around a connection without changing the classification target: the
decision remains edge-level. It also separates information attached to the communicating
endpoints from information carried by the flow itself.

E-GraphSAGE is an important reference point for this formulation. It applies a GNN to
intrusion detection and classifies an edge from the concatenation of the embeddings of its
two endpoints \citep{lo2022egraphsage}. The edge's own features do not reach the
classifier. That choice gives the classification head a purely endpoint-derived view of the
connection. A direct consequence is that two parallel edges with the same ordered endpoints
present the same pair of endpoint embeddings to the edge classifier, even if the edges
differ in attributes that are not otherwise represented at the head. The architecture can
express structural context around an endpoint pair, but its final edge representation does
not preserve an explicit channel for the attributes of the particular edge being classified.

TE-G-SAGE moves to an edge-aware GraphSAGE formulation: its supervised edge head consumes
the edge's own features, and the model is paired with SHAP-based explanation, an attribution
method that divides a prediction into per-feature contributions \citep{tegsage2025}. This changes the representational question. Once the edge itself is
allowed to carry predictive information, the model can distinguish endpoint context from
flow-specific attributes rather than requiring the classifier to infer the latter
indirectly from node embeddings.

For network-flow intrusion detection that distinction matters, because the graph and the
individual connection encode different kinds of evidence. The graph describes who
communicates with whom and under what relational context. The edge describes the flow being
classified. The graph-based IDS literature therefore motivates an edge-level representation
retaining both endpoint structure and edge-specific content. These two architectures are also
the ones re-trained on our own aggregated representation as baselines in
Section~\ref{sec:res-baselines}. The next question is whether edge content has to remain
numeric, or whether some of it can be represented semantically.

\subsection{Language models over security text and network traffic}

Cybersecurity language encoders provide one route to domain-specific semantics. CySecBERT
is a BERT-class encoder adapted to cybersecurity text \citep{bayer2024cysecbert}. Its
relevance here is representational: a pretrained cybersecurity language encoder maps
controlled security-related text into a dense embedding without requiring the graph encoder
to interpret textual tokens itself.

A separate line of work treats traffic as a language-like sequence. BSTFNet applies a
BERT-class model, ET-BERT, to serialised packet-level traffic and combines that
representation with a spatiotemporal branch for encrypted-traffic classification
\citep{huang2024bstfnet}. It uses no graph encoder. The precedent that matters is the
serialisation step itself: network observations that begin as structured traffic fields can
be converted into a controlled token sequence and processed by a BERT-class encoder.

These two lines solve different problems. Domain adaptation addresses the vocabulary and
semantics of cybersecurity text. Traffic serialisation creates text-like input from network
measurements. Together they support a semantic branch that encodes
a flow description as text rather than treating every field as a scalar magnitude. They do not, by
themselves, determine how that semantic representation should interact with graph structure.

That interaction is the central fusion problem. Once a system holds one representation
derived from topology and another derived from text, the design space is no longer only
about which encoder to choose. It is about where the two modalities meet, and whether their
relative influence is fixed across samples or allowed to change from one sample to the next.

\subsection{Fusing structural and semantic representations}

Fusion methods organise along two related axes: how tightly the modalities interact, and
whether the fusion rule assigns the same modality weighting to every sample or adapts that
weighting to the input. The second axis is the more important one here. Four recurring
families cover the relevant design space.

The simplest family fixes the interface between the branches in advance. BertGCN is the
clearest example: its final prediction linearly interpolates a GCN prediction and a BERT
prediction using a single global hyperparameter \citep{lin2021bertgcn}. Every sample receives
the same interpolation coefficient. Such designs keep the encoders modular, but the fusion
rule never decides that one modality should matter more for one input and less for another.

A second family introduces input-dependent gates. The Gated Multimodal Unit uses
multiplicative gates to determine how each modality contributes to the fused activation
\citep{arevalo2017gmu}. This is the canonical primitive for sample-adaptive modality
weighting. A graph-and-text instance appears in the Dual-Stream Graph-Augmented Transformer
for fake-news detection, where BERT runs in parallel with a GAT and Graph Transformer branch,
and a trainable scoring function emits a scalar weight for each modality for each input
\citep{moorthy2025dualstream}. Mechanically this is the closest analogue in the literature
surveyed here to a fusion stage whose balance changes per sample.

A third family uses attention to let the modality representations interact directly.
RAGFormer stacks the two branch embeddings as a short token sequence and applies
self-attention over them with a residual connection \citep{li2024ragformer}. CAST uses
cross-attention between graph and text representations in materials prediction and motivates
the text branch by arguing that graph conversion can discard global information which
textual descriptions restore \citep{lee2025cast}. The semantic representation is therefore
not a duplicate encoding of graph structure. BiGTex goes further and interleaves a two-way
exchange: structural information enters the language model as a soft prompt, and text
returns to the graph side through cross-attention with the query supplied from the graph
side \citep{beiranvand2026bigtex}. GMLM runs the exchange in both directions at once: the
language embeddings attend to the graph embeddings and the graph embeddings attend to the
language embeddings, and the two attended results are concatenated before classification, a
design its authors set explicitly against simple concatenation \citep{sinha2025gmlm}.
Attention-based interaction differs from global interpolation because its interaction weights
are derived from the representations of the current input.

The fourth family removes a distinct late-fusion stage altogether by interleaving structural
and semantic processing throughout the encoder. DET uses exactly the terms ``structural
encoder'' and ``semantic encoder'' and argues against late fusion of separate encoders,
preferring interleaved interaction \citep{guo2024det}. GL-Fusion sits at the opposite end of
the coupling spectrum from independent towers: message passing is folded into every
transformer layer, leaving no separate fusion module \citep{yang2024glfusion}. These designs
mix the modalities repeatedly instead of preserving two independent representations until
classification.

The taxonomy separates two questions that are often collapsed. One is where graph and text
meet: after separate encoding, inside a dedicated attention mechanism, or repeatedly
throughout the encoder. The other is whether the relative influence of the modalities is
globally fixed or varies with the sample. The fusion stage built in this report belongs to the
second family. A gate reads both projected representations and emits a pair of weights that sum
to one, so the weighting between structure and semantics is decided separately for every edge
rather than by a single coefficient shared across the dataset. The second distinction becomes
sharper on moving
from general graph--text systems to intrusion detection, where the available GNN--language
combinations tend to assign the two components sequential roles rather than fusing their
representations.

\subsection{GNN and LLM combinations in intrusion detection}

Within the intrusion-detection works surveyed here, graph models and language models are
combined sequentially rather than fused as two representations of the same edge-level
decision. XG-NID is the clearest example. Its heterogeneous GNN performs the intrusion
classification, and an LLM is used only after prediction, to produce human-readable
explanations \citep{farrukh2025xgnid}. The system's ``dual modality'' refers to flow-level
and packet-level data inside the graph pipeline, not to graph and text embeddings fused for
classification.

A poster by \citet{zhan2025agentgnn}, which describes its aim as enhancing GNN ``robustness'' for network intrusion detection, reverses the order. % numcheck: allow -- "robustness" is quoted from that poster's own title, not a claim about our system
An LLM agent analyses textual descriptions of node interactions and flags, removes, or
reweights suspicious nodes and edges before the GNN processes the graph. The language component
acts as a preprocessing filter: its output changes the graph presented to the GNN, but the
two components are never combined as concurrent representations at the classifier.

Other intrusion-detection systems occupy adjacent points in the design space without forming
a GNN--language pair at all. GCN-2-Former combines a GCN with a temporal transformer and
fuses the two by global average pooling and concatenation \citep{gcn2former2025}. The second
tower is a temporal transformer, not a language model. CPS-IDS uses gated fusion between a
DeBERTa payload-semantic branch and a statistical branch \citep{xiong2026cpsids}, so it has
both a language-based semantic component and an explicit fusion mechanism, but no graph
encoder.

The distinction is functional rather than terminological. In this literature the language
model either explains a GNN output after classification or alters graph inputs before
message passing. The parallel-fusion examples occupy nearby configurations but either
replace the language model with a non-language temporal transformer or replace the graph
encoder with a statistical branch. Representation-level fusion and graph--language feedback
therefore have to be examined in the broader graph-learning literature rather than inferred
from the intrusion-detection examples alone.

\subsection{Iterative GNN--LLM feedback}

The broader GNN--LLM literature provides four relevant feedback approaches.

DAS is the closest prior work to the feedback mechanism used in this report. It implements
an explicit closed feedback loop between a frozen GNN and an LLM: GNN predictions provide
implicit supervision that steers semantic refinement, and the revised semantics then re-enter
the same graph learner \citep{thapaliya2025das}. DAS also writes each node's structural
statistics out as a sentence, using degree, betweenness, closeness, clustering coefficient and
square clustering coefficient, and gives a Majorization--Minimization argument, the technique of
repeatedly minimising an easier upper bound in place of a hard objective, under which repeated
refinement never increases the surrogate objective it defines for the task. Its setting is
node-level, and it is evaluated on graphs that carry node text and on graphs that do not, where
that structural sentence is what supplies the missing text. Its reported gains concentrate on
the structure-dominated graphs.

LOGIN already couples uncertainty selection to language-model consultation. It uses
MC-dropout variance, the spread of predictions obtained by running the network several times
with dropout left switched on, to identify uncertain nodes. It then prompts the LLM with the
node's own text and a description of its two-hop neighbourhood, and receives a predicted label
together with a free-text explanation
\citep{qiao2025login}. Incorporation depends on correctness: if the LLM's label matches the
ground-truth label, the explanation is appended to the node text and the augmented text is
re-embedded. If the label is wrong, edges in the node's one-hop neighbourhood are pruned.
The LLM remains outside autograd and edits the learning inputs between training rounds.


GLANCE makes the decision to consult the language model itself learnable. A lightweight
router decides per node whether querying is worthwhile from inexpensive node features, and
an advantage-based objective trains that router on the outcome of querying versus not
querying, because the LLM call is non-differentiable \citep{loveland2025glance}. Its gains
concentrate on heterophilous and structurally ambiguous nodes. This addresses the same
allocation problem as uncertainty-based consultation, directing language-model effort at
selected graph elements rather than applying it uniformly, with a learned criterion in
place of a fixed uncertainty rule.

RoGRAD places iterative refinement, rather than static one-shot signal injection, at the
centre of its formulation \citep{wang2025rograd}. It studies that refinement under graph
deficiencies including noise, sparsity, and label scarcity. Its abstract claims priority
for moving language-model augmentation of graphs from static injection to an iterative
refinement paradigm.

\subsection{Positioning}

The contribution is an edge-level intrusion-detection instantiation, and a measurement
question about how such a feedback channel is used. The comparison has to be made
at the level of configuration, not of individual mechanisms.

The four feedback anchors operate at node level, whereas the present setting treats a
network flow as an edge and makes the classification decision there. Their semantic
consultants are prompted generative LLMs. Here the semantic branch is a frozen sentence
encoder, and the consultant reading its embeddings is a classification head trained on them.
An earlier configuration scored those embeddings against class prototypes instead, and both
are reported. Their feedback actions refine
semantic inputs or edit graph inputs. Here the accepted semantic signal modifies the edge
representations consumed by the next graph update. These are differences of instantiation,
not claims that the underlying ideas originated here
\citep{qiao2025login,thapaliya2025das,loveland2025glance,wang2025rograd}.

Two signals travel in opposite directions in the resulting loop. The graph model's
per-edge uncertainty determines \emph{which} edges are
routed to the semantic branch, and that selection is recomputed as predictions shift. The
semantic branch returns a per-edge judgement that re-enters the graph model's edge
representations. This is the routing-level bidirectionality defined in
Section~\ref{sec:introduction}, and the content-level limit stated there applies here as well:
the consultant's verdict on an edge depends on that edge's text alone.

The final classification stage adds a further dimension to the configuration. Structural and
semantic edge representations are combined through gated fusion with
sample-dependent modality weighting, rather than assigning the semantic component a
preprocessing or post-hoc role. Sample-adaptive fusion is itself established outside this
setting: gated multimodal units provide input-dependent modality control, and the
Dual-Stream Graph-Augmented Transformer assigns per-input scalar modality weights to
parallel text and graph representations \citep{arevalo2017gmu,moorthy2025dualstream}.
Representation-level graph--text interaction is likewise established, through
self-attention, cross-attention, and interleaved exchange
\citep{li2024ragformer,beiranvand2026bigtex,lee2025cast}. None of those mechanisms is a
contribution claim here.

Within the limits of the literature search represented by this survey, we are not aware of
prior work combining all three of the following in one system: bidirectional iterative
coupling between a graph model and a semantic branch, sample-adaptive gated fusion for the
final classification, and edge-level network intrusion detection. This is a statement about
the configuration space our search covered. It does not imply a performance advantage, and
it does not change the prior-art status of any individual component.

That positioning leaves a narrower research question. Existing feedback work asks whether a
language-model-driven refinement process improves the graph-learning task, with the model
revising semantics, modifying inputs, or selectively augmenting graph elements
\citep{qiao2025login,thapaliya2025das,loveland2025glance,wang2025rograd}. The present work
instead instantiates that loop in network-flow intrusion detection with a deliberately more
constrained consultant, and asks a diagnostic question: whether such a consultant exploits a
feedback channel that has the capacity to influence subsequent graph updates. That question
is underexplored in the work surveyed here. No claim about its answer belongs in this
section.

The contribution is best described as instantiation plus measurement. The system instantiates
bidirectional feedback and sample-adaptive structural--semantic fusion for edge-level network
intrusion detection, then measures whether the constrained semantic consultant makes use of
the channel it is given. It does not rest on inventing GNN--LLM feedback loops, iterative
semantic refinement, natural-language descriptions of graph statistics, or selective semantic
consultation. Those mechanisms have clear precedents. The point of this configuration is to
make their interaction measurable in an edge-level network-IDS setting.
